% !Mode:: "TeX:UTF-8"

%%% 此部分需要自行填写: (1) 中文摘要及关键词 (2) 英文摘要及关键词
%%%%%%%%%%%%%%%%%%%%%%%%%%%%%
%%% -------------  英文封面 (无需改动)-------------   %%%
%%%%%%%%%%%%%%%%%%%%%%%%%%%%%
\thispagestyle{empty}
\renewcommand{\baselinestretch}{1.5}  %下文的行距
\vspace*{0.5cm}
\begin{center}
{\Large \bf BACHELOR'S DEGREE THESIS \\[1ex] OF WUHAN UNIVERSITY }
\end{center}
\vspace{2.5cm}
\begin{center}{\zihao{2} \the\Etitle \par}\end{center}

\vfill

\begin{center}
\zihao{4}
\begin{tabular}{ r l }
 School (Department): & {\sc \the\Eschoolname}\\
  Major:          &   {\sc\the\Emajor}  \\
 Candidate:      &  {\sc \the\Eauthor}      \\
 Supervisor:     &  {\sc \the\Esupervisor}
\end{tabular}

\vspace*{2cm}
\begin{center}
   \ifprint % 文档打印, 使用黑白校徽.
  \includegraphics[height=4cm]{whu.eps}       %%  黑白的.
  \else
  \includegraphics[height=4cm]{whulogo.eps} %%  彩色的.
  \fi
\end{center}


\zihao{-2}
%\the\Schoolname\\
{\sc Wuhan University}

\vspace*{1.0cm}

\the\Edate

\end{center}
%%% 郑重声明部分无需改动

%%%---- 郑重声明 (无需改动)------------------------------------%
\newpage
\vspace*{20pt}
\begin{center}{\ziju{0.8}\textbf{\songti\zihao{2} 郑重声明}}\end{center}
\par\vspace*{30pt}
\renewcommand{\baselinestretch}{2}

{\zihao{4}%

本人呈交的学位论文, 是在导师的指导下, 独立进行研究工作所取得的成果,
所有数据、图片资料真实可靠. 尽我所知, 除文中已经注明引用的内容外,
本学位论文的研究成果不包含他人享有著作权的内容.
对本论文所涉及的研究工作做出贡献的其他个人和集体,
均已在文中以明确的方式标明. 本学位论文的知识产权归属于培养单位.\\[2cm]

\hspace*{1cm}本人签名: $\underline{\hspace{3.5cm}}$
\hspace{2cm}日期: $\underline{\hspace{3.5cm}}$\hfill\par}
%------------------------------------------------------------------------------
\baselineskip=23pt  % 正文行距为 23 磅
%------------------------------------------------------------------------------





%%======中文摘要===========================%
\begin{cnabstract}
智能终端设备的盛行使得移动平台的应用



\end{cnabstract}
\par
\vspace*{2em}


%%%%--  关键词 -----------------------------------------%%%%%%%%
%%%%-- 注意: 每个关键词之间用“；”分开，最后一个关键词不打标点符号
\cnkeywords{毕业论文; \LaTeX{}; 模板;  }


%%====英文摘要==========================%


\begin{enabstract}
The prevalence of mobile platforms, the large market
share of Android, plus the openness of the Android Market makes it a hot target for malware attacks. Once a malware sample has been identified, it is critical to quickly
reveal its malicious intent and inner workings. 

In this
paper we present DroidScope, an Android analysis platform that continues the tradition of virtualization-based
malware analysis. Unlike current desktop malware analysis platforms, DroidScope reconstructs both the OSlevel and Java-level semantics simultaneously and seamlessly. To facilitate custom analysis, DroidScope exports three tiered APIs that mirror the three levels of an
Android device: hardware, OS and Dalvik Virtual Machine. On top of DroidScope, we further developed several analysis tools to collect detailed native and Dalvik
instruction traces, profile API-level activity, and track information leakage through both the Java and native components using taint analysis. These tools have proven to
be effective in analyzing real world malware samples and
incur reasonably low performance overheads..

\end{enabstract}
\par
\vspace*{2em}

%%%%%-- Key words --------------------------------------%%%%%%%
%%%%-- 注意: 每个关键词之间用“；”分开，最后一个关键词不打标点符号
 \enkeywords{\LaTeX{};  }
